\chapter{Getting Started}
\label{ch:getting_started}

Welcome to the IISER Kolkata MS and PhD thesis template. This document serves as the official documentation, demonstrating how to use the template effectively to produce a beautiful, correctly formatted thesis.

\section{Project Structure}
The template is divided into several modular components to keep your workspace clean:
\begin{itemize}
    \item \mintinline{text}{main.tex}: The root document where you configure metadata and include other files.
    \item \mintinline{text}{IISERKthesis.cls}: The core LaTeX class file that manages page layout, margins, headers, and frontmatter.
    \item \mintinline{text}{IISERKphysics.sty}: A custom package that loads mathematics tools, defines theorem environments, and configures code blocks.
    \item \mintinline{text}{frontmatter/}: Contains individual files for your abstract, acknowledgements, certificate, and declaration.
    \item \mintinline{text}{chapters/} and \mintinline{text}{appendices/}: Directories to house your main textual content.
    \item \mintinline{text}{references.bib}: BibLaTeX database containing your literature citations.
\end{itemize}

\section{Metadata Configuration}
In \mintinline{text}{main.tex}, you will find several macros to define your personal details. These automatically populate the title page, declaration, and certificate:
\begin{minted}{latex}
\title{Your Thesis Title Here}
\author{Your Name}
\rollno{12MS345}
\program{5th Year BS-MS}
\department{Department of Physical Sciences}
\supervisor{Prof. Supervisor Name}
\end{minted}

\section{Compilation}
This template relies on \mintinline{text}{biber} for bibliography management and \mintinline{latex}{minted} for syntax-highlighted code blocks. Because of \mintinline{latex}{minted}, you \emph{must} compile with the \mintinline{bash}{-shell-escape} flag.

The easiest way to compile is using the included \mintinline{text}{.latexmkrc} configuration:
\begin{minted}{bash}
latexmk main.tex
\end{minted}
This automates multiple compilation passes and stores all auxiliary files (like \mintinline{text}{.aux} and \mintinline{text}{.log}) neatly in a hidden \mintinline{text}{.build_tex/} folder.

\section{Class Options and Toggles}
\label{sec:class_options}
The template automatically generates several lists at the beginning of the document (Table of Contents, List of Figures, List of Tables, List of Listings, Nomenclature, and a Todo list).

You can easily disable any of these by passing toggles to the document class in \mintinline{text}{main.tex}:
\begin{minted}{latex}
% Disables the Todo list and List of Tables
\documentclass[notodo, nolot]{IISERKthesis}
\end{minted}

Available toggles include:
\begin{itemize}
    \item \mintinline{latex}{notodo}: Completely disables the \mintinline{latex}{todonotes} package and removes the Todo List. Ensure you use this flag for your final submission!
    \item \mintinline{latex}{nolof}: Removes the List of Figures.
    \item \mintinline{latex}{nolot}: Removes the List of Tables.
    \item \mintinline{latex}{nolol}: Removes the List of Listings.
    \item \mintinline{latex}{nonom}: Removes the Nomenclature (List of Symbols).
\end{itemize}
